% Generated draft: requires mathematical and source review.
\documentclass[11pt]{article}
\usepackage[margin=25mm]{geometry}
\usepackage{fontspec}
\setmainfont{Latin Modern Roman}
\usepackage{amsmath,amssymb,mathtools,unicode-math}
\setmathfont{Latin Modern Math}
\usepackage{enumitem,needspace,xurl,hyperref}
\setlength{\parindent}{0pt}
\setlength{\parskip}{5pt}
\setlength{\emergencystretch}{3em}
\newcommand{\N}{\mathbb{N}}
\newcommand{\Z}{\mathbb{Z}}
\newcommand{\Q}{\mathbb{Q}}
\newcommand{\R}{\mathbb{R}}
\newcommand{\C}{\mathbb{C}}
\newcommand{\F}{\mathbb{F}}
\newcommand{\A}{\mathbb{A}}
\newcommand{\PP}{\mathbb{P}}
\newcommand{\E}{\mathbb{E}}
\newcommand{\eps}{\varepsilon}
\newcommand{\dd}{\,\mathrm{d}}
\begin{document}
\section*{UnsolvedMath ID: 700660}
\textbf{Source problem code:} catalogue entry
\subsection{Definitions and mathematical statement}
Let \(R\) be an associative ring with identity, and let \(n\in\Z\) satisfy \(n\geq 3\). Write \(GL_n(R)\) for the group of invertible \(n\times n\) matrices over \(R\). For distinct indices \(i,j\in\{1,\ldots,n\}\), let \(e_{ij}\) be the corresponding matrix unit, and define the elementary subgroup by
\[
E_n(R)=\left\langle I_n+r e_{ij}:r\in R,\ i\ne j\right\rangle.
\]
For subgroups \(A,B\leq GL_n(R)\), set
\[
[A,B]=\left\langle a^{-1}b^{-1}ab:a\in A,\ b\in B\right\rangle.
\]
A subgroup \(N\leq GL_n(R)\) is normalized by \(E_n(R)\) if \(gNg^{-1}=N\) for every \(g\in E_n(R)\).

The supplied problem, proposed by V. M. Petechuk, asks whether, for every such \(R,n,N\),
\[
E_n(R)\leq N
\quad\text{or}\quad
\exists\, I\lneq R\text{ an ideal such that }
\Lambda_I[N,E_n(R)]=1.
\]
The supplied record does not define \(\Lambda_I\), specify its action on a subgroup, or explicitly specify whether ideals are two-sided. Consequently, the second alternative is not fully determined by the supplied evidence.

If \(\Lambda_I\) denotes the homomorphism induced by reduction modulo a two-sided ideal \(I\), define
\[
\pi_I:GL_n(R)\longrightarrow GL_n(R/I),\qquad
(a_{ij})\longmapsto(a_{ij}+I).
\]
Under that interpretation, the requested assertion is precisely
\[
\forall R\ \forall n\in\Z_{\geq3}\ \forall N\leq GL_n(R),\qquad
\left(\forall g\in E_n(R),\ gNg^{-1}=N\right)
\Longrightarrow
\left(
E_n(R)\leq N\ \lor\
\exists\, I\lneq R\text{ a two-sided ideal},\
\pi_I\bigl([N,E_n(R)]\bigr)=\{I_n\}
\right).
\]
Here \(R\) ranges over associative rings with identity, and the identity matrix in the last expression is over \(R/I\). This reduction-map formulation is conditional on confirmation of the undefined source notation; it is not an established identification of \(\Lambda_I\).
\subsection{Short English statement}
Petechuk asks whether every subgroup of \(GL_n(R)\), for an associative unital ring \(R\) and \(n\geq3\), that is normalized by the elementary subgroup either contains that subgroup or satisfies the stated commutator condition for some proper ideal. The source leaves \(\Lambda_I\) undefined; a precise formulation using reduction modulo a two-sided ideal is given conditionally and requires source confirmation.
\subsection{Sources}
\begin{itemize}
\item Supplied local UnsolvedMath registry record, permanent ID \(700660\), exact source problem code \texttt{catalogue entry}, title \emph{Kourovka Notebook Problem 10.45}; upstream revision \texttt{d371ee917ac9647fbb06cd6fc529deec6f3dda82}.
\item \emph{The Kourovka Notebook (Unsolved Problems in Group Theory)}, original problem 10.45, 10th issue (1986); proposed by V. M. Petechuk. The supplied background identifies the 21st issue edition (2026), \texttt{arXiv:1401.0300v48}. The supplied collection description names E. I. Khukhro and V. D. Mazurov as editors.
\item The registry labels this problem \emph{open} and records an update on 2026-10-16. This is supplied registry information, not an independently verified mathematical or literature assessment. No source was independently consulted for this draft.
\end{itemize}
\end{document}
